\subsection{消退期}

射精过后，男性的身体进入消退期，逐渐回到未唤起前的状态。玛斯特斯和约翰逊夫妇把男性的消退分成两个阶段：第一阶段，阴茎仍然保持半勃起，这种状态可以持续一段时间，此时若继续给予温和的刺激，一部分男性还能较快地重新勃起；到了第二阶段，勃起会迅速而完全地消退，充血散去，睾丸降回原位，阴囊的皮肤重新松弛下来。

伴随消退期而来的是男性特有的不应期——射精后的一段时间里，无论受到多大的性刺激，都无法再次唤起勃起与射精。这与射精后催乳素升高、神经系统从兴奋转向抑制有关。不应期长短因人、因龄而异：年轻人可能只有几分钟到几十分钟，随年龄增长可延长至数小时甚至更久。这是正常的生理差异，不是"行不行"的指标。不应期的原理与"二次性爱"的管理，本书前文已设有专节（"多巴胺、催乳素与性后不应期"与"二次性爱与不应期的管理"两节），此处不再展开。

需要说明的是，男性只有在尚未射精时才能退回平台期、让性紧张反复徘徊；射精一旦发生，消退期就不可避免。消退期里的困倦与放松感非常普遍——这正是高潮后催乳素升高的表现之一。与女性相比，男性的消退要"彻底"得多：女性在高潮后若继续受到刺激，可以退回平台期、再度达到高潮，而男性必须先度过不应期。因此性生活中常常出现这样的节奏差：男性高潮后想睡，女性却还期待温存与爱抚。懂得这个差异，把事后的时间留给拥抱与交谈，对双方都是一种体贴。
